\documentclass[11pt]{article}

\usepackage[margin=0.72in]{geometry}
\usepackage{amsmath,amssymb,amsthm}
\usepackage{hyperref}
\usepackage{booktabs}

\hypersetup{colorlinks=true,linkcolor=blue,citecolor=blue,urlcolor=blue}

\newtheorem{theorem}{Theorem}
\newtheorem{proposition}[theorem]{Proposition}
\newtheorem{corollary}[theorem]{Corollary}

\title{The Hadamard Sign Kernel of the Tetrahedral Frame}
\author{Stephen Nelson\\
\small Field Dynamics Research\\
\small\texttt{stephen@fielddynamics.org}}
\date{}

\begin{document}
\maketitle

\begin{abstract}
The $4\times 4$ Hadamard matrix $H_4$ appears independently in
three literatures: as the parity-restricted Walsh-Hadamard kernel
in signal processing, as the sign skeleton of the isotropic tetrad
(vierbein) in general relativity, and as the axis structure of the
stellated octahedron in discrete geometry.  This note observes that
these three appearances are the same object.  The full
Walsh-Hadamard transform on the cube $\mathbb Z_2^3$ is the
$8\times 8$ matrix $H_8$.  Restricting to one tetrahedral parity
sector ($xyz=+1$) collapses the eight Walsh characters into four,
and the resulting evaluation matrix is $H_4$.  The Sylvester
factorisation $H_8=H_2\otimes H_4$ encodes this decomposition.
The three non-constant rows of $H_4$ are the spatial symmetry axes
of the stellated octahedron and the face normals of the isotropic
tetrahedral frame forced by spatial isotropy in $D=3$ dimensions.
The identification requires no free parameter and holds exclusively
at the frame dimension $K=D+1=4$.
\end{abstract}


\section{The cube, the transform, and the compound}

The Walsh-Hadamard transform (WHT) is the discrete Fourier
transform of the elementary abelian 2-group $\mathbb Z_2^n$.  At
$n=3$, its domain is the cube: the eight vertices
$\{(\pm1,\pm1,\pm1)\}$, or equivalently the 3-bit binary strings
$\mathbb F_2^3$.  The Walsh functions are the group characters
$\chi_S(x) = \prod_{i\in S} x_i$, indexed by subsets
$S\subseteq\{1,2,3\}$.  The transform and the cube are
inseparable; the WHT is defined on $\mathbb Z_2^3$ and nowhere
else at $n=3$~\cite{Walsh,Sylvester}.  The full transform matrix
at $n=3$ is $H_8$, an $8\times 8$ matrix acting on all eight cube
vertices.  The $4\times 4$ matrix $H_4$ that is the subject of
this note arises not as the full cube transform but as its
restriction to one parity sector, as shown in Section~3.

The same cube contains the stellated octahedron.  The parity
function $p(x,y,z) = xyz$ splits the eight vertices into two sets
of four:
\begin{align*}
\text{TetA}\ (p=+1)&\colon\
(+,+,+),\ (+,-,-),\ (-,+,-),\ (-,-,+),\\
\text{TetB}\ (p=-1)&\colon\
(-,-,-),\ (-,+,+),\ (+,-,+),\ (+,+,-).
\end{align*}
Any two vertices within the same set differ in exactly two
coordinates, giving pairwise inner product $-1$.  Four points in
$\mathbb R^3$ with constant pairwise inner product form a regular
tetrahedron~\cite{Coxeter}.  The compound of TetA and TetB is the
stellated octahedron (stella octangula), inscribed in the cube on
the same eight vertices.  This is classical discrete
geometry~\cite{Kepler,Coxeter}.

The observation of this note is that the symmetry axes of this
compound, already present inside the cube, are the rows of the
$4\times 4$ Hadamard matrix.


\section{Balanced partitions and the Hadamard matrix}

A balanced partition of a set of four objects into two pairs
admits exactly three distinct configurations:
\[
\{1,2\}\mid\{3,4\},\qquad
\{1,3\}\mid\{2,4\},\qquad
\{1,4\}\mid\{2,3\}.
\]
This count is $\binom{4}{2}/2=3$.  Each partition can be encoded
as a $\pm1$ vector of length four, assigning $+1$ to one pair and
$-1$ to the other.  Together with the constant vector
$(+1,+1,+1,+1)$, these four vectors form the rows of the
$4\times 4$ Hadamard matrix:
\[
H_4 = \begin{pmatrix}
+1 & +1 & +1 & +1 \\
+1 & +1 & -1 & -1 \\
+1 & -1 & +1 & -1 \\
+1 & -1 & -1 & +1
\end{pmatrix}.
\]
Orthogonality is immediate: the constant row is orthogonal to
every balanced row (equal numbers of $+1$ and $-1$), and any two
balanced rows agree on exactly two entries and disagree on two,
giving inner product zero.  Therefore $H_4 H_4^T = 4I_4$.

The three balanced rows are the three independent ways to split
four objects into equal halves.  There is no fourth; the next
partition $\{1,2,3\}\mid\{4\}$ is unbalanced and not orthogonal
to the constant row.


\section{The axes of the stellated octahedron}

The stellated octahedron has three mutually perpendicular symmetry
axes, connecting opposite edge-midpoints of the cube.  Each axis
separates the four TetA vertices (equivalently, the four TetB
vertices) into two pairs: those with $+1$ along that axis and
those with $-1$.

Explicitly, the four TetA face normals are:
\[
u_0 = (+,+,+),\quad
u_1 = (+,-,-),\quad
u_2 = (-,+,-),\quad
u_3 = (-,-,+).
\]
Reading the sign of each normal along the three spatial axes
gives, for each axis, a balanced partition of the four faces into
two pairs:
\begin{center}
\begin{tabular}{lccccc}
\toprule
& $u_0$ & $u_1$ & $u_2$ & $u_3$ & Partition \\
\midrule
$x$-axis (sign of first component)
  & $+$ & $+$ & $-$ & $-$ & $\{0,1\}\mid\{2,3\}$ \\
$y$-axis (sign of second component)
  & $+$ & $-$ & $+$ & $-$ & $\{0,2\}\mid\{1,3\}$ \\
$z$-axis (sign of third component)
  & $+$ & $-$ & $-$ & $+$ & $\{0,3\}\mid\{1,2\}$ \\
\bottomrule
\end{tabular}
\end{center}
These are the three balanced partitions of Section~2.  Adding the
constant row (all faces positive) recovers $H_4$ exactly.

\begin{theorem}
The three spatial symmetry axes of the stellated octahedron,
expressed as sign patterns across the four tetrahedral face
normals, together with the constant temporal axis, form the rows
of $H_4$.
\end{theorem}

\begin{proof}
The table above exhibits the identification.  Orthogonality
follows from the balanced-partition structure of Section~2.
\end{proof}

The geometric identification has an algebraic counterpart that
explains why $H_4$ appears inside the $H_8$ Walsh domain of the
full cube.

\begin{theorem}[Parity-restricted Walsh kernel]
Let $T_+=\{(x,y,z)\in\{\pm1\}^3:xyz=+1\}$.  The restriction of
the eight Walsh characters on $\mathbb Z_2^3$ to $T_+$ identifies
complementary character pairs:
\[
        1\sim xyz,\qquad x\sim yz,\qquad y\sim xz,\qquad z\sim xy.
\]
The restricted character algebra therefore has four independent
modes $\{1,x,y,z\}$.  Evaluating these four modes on the four
vertices of $T_+$ gives $H_4$.
\end{theorem}

\begin{proof}
On $T_+$, $xyz=+1$ by definition.  Multiplying both sides by any
single variable: $yz = x$, $xz = y$, $xy = z$.  So
$\chi_{\{i,j\}}|_{T_+} = \chi_{\{k\}}|_{T_+}$ for
$\{i,j,k\}=\{1,2,3\}$, and
$\chi_{\{1,2,3\}}|_{T_+} = \chi_\emptyset|_{T_+} = 1$.  The
eight Walsh characters collapse to four.  Evaluating
$\{1,x,y,z\}$ on the four TetA vertices $(+,+,+)$, $(+,-,-)$,
$(-,+,-)$, $(-,-,+)$ gives the rows of $H_4$ by direct
computation.
\end{proof}

After reordering the row and column bases so that the first factor
is the parity character $xyz$, the Sylvester factorisation
\[
        H_8 \cong H_2\otimes H_4
\]
encodes this decomposition.  The $H_2$ factor carries the
TetA/TetB grading ($xyz=\pm 1$), while the $H_4$ factor carries
the frame structure within each parity sector.  This is the
stellated octahedron's self-similar structure expressed as a tensor
product: the same cube, the same three dimensions, decomposed into
parity grading times frame read.


\section{The isotropic tetrad frame}

In the standard tetrad formulation of general relativity, the
gravitational field is described by a frame field $e^a{}_\mu$ with
$K=D+1$ internal legs~\cite{MTW}.  In $D=3$ spatial dimensions,
$K=4$.

The Lorentz signature $(-,+,+,+)$ maps null directions to the
unit sphere $S^2$~\cite{NelsonA}.  Four unit directions on $S^2$
that introduce no net preferred spatial direction and treat all
directions equivalently are uniquely (up to $O(3)$) the vertices
of a regular tetrahedron, with pairwise inner product
$c=-1/3$~\cite{NelsonA,Coxeter}.  The face normals of this
isotropic tetrahedral frame are, up to an overall normalisation
factor $1/\sqrt{3}$, the four vectors $u_0,\ldots,u_3$ of the
previous section.

The sign pattern of the spatial tetrad components $e^I{}_j$ across
the four frame legs $I\in\{0,1,2,3\}$ and three spatial directions
$j\in\{1,2,3\}$ is therefore:
\[
\operatorname{sgn}(e^I{}_j) =
\begin{pmatrix}
+ & + & + \\
+ & - & - \\
- & + & - \\
- & - & +
\end{pmatrix}
= \text{spatial block of } H_4.
\]
The temporal leg $e^0{}_0=1$ (with $e^0{}_j=0$) supplies the
constant row.  The sign skeleton of the isotropic tetrahedral
frame is $H_4$, up to the normalisation factor $1/\sqrt{3}$ on
the spatial block.  The geometric frame is
$e^I{}_j = u_I^j = (\pm1,\pm1,\pm1)/\sqrt{3}$; its sign kernel
is $H_4$.

\begin{corollary}
The parity-restricted Walsh-Hadamard kernel has the same sign
structure as the change of basis from coordinate axes to the
isotropic tetrahedral frame.  The FWHT kernel and the tetrahedral
tetrad frame therefore share the same Hadamard sign skeleton, up
to ordering, transpose, and normalisation.
\end{corollary}


\section{Uniqueness at $D=3$}

The identification rests on three coincidences that hold
simultaneously only at $D=3$.

\textbf{The simplex has $K=D+1$ faces.}  At $D=3$, $K=4$.  The
four faces admit exactly $\binom{4}{2}/2=3=D$ independent balanced
partitions, matching the spatial dimension.  At $K=3$ ($D=2$):
$\binom{3}{2}/2=3/2$, not an integer.  At $K=5$ ($D=4$):
$\binom{5}{2}/2=5\neq 4=D$.  Only at $K=4$ does the
balanced-partition count equal the spatial dimension.

\textbf{The cube dimension matches the simplex face count.}  The
cube $\mathbb Z_2^D$ has $2^D$ vertices.  Two tetrahedra have
$2K=2(D+1)$ vertices.  These are equal only when
$2(D+1)=2^D$, which holds at $D=3$ and no other positive
integer~\cite{NelsonA}.

\textbf{The Hadamard matrix exists at order $K$.}  A necessary
condition for a real $\pm1$ Hadamard matrix of order $K>2$ is
$K\equiv 0\pmod{4}$~\cite{Hadamard}.  At $K=4$ this is satisfied.
At $K=3$ or $K=5$ no such matrix exists.

\begin{proposition}
$K=4$ is the unique frame dimension at which the stellated
octahedron axes form a Hadamard matrix matching the spatial
dimension.
\end{proposition}


\section{Discussion}

The content of this note is an observation, not a construction.
The Walsh-Hadamard transform has been used in signal processing
since Walsh (1923) and Hadamard (1893).  The stellated octahedron
has been studied since Kepler (1611).  The isotropic tetrahedral
frame is a standard object in tetrad
gravity~\cite{MTW,Hehl1976}.  Each appears in a mature literature
with extensive bibliography.  What has not appeared, to the
author's knowledge, is the identification of these three objects as
the same $4\times 4$ matrix.

The identification has a structural consequence.  The WHT is the
native Fourier analysis of the cube $\mathbb Z_2^3$.  The
stellated octahedron is the parity structure of the same cube.
The tetrad frame is the gravitational frame on the same four
directions.  All three inhabit the same eight-vertex geometry,
and $H_4$ is their shared algebraic expression.

The parity-restricted kernel has four modes: three spatial
characters $\{x,y,z\}$ and one constant scalar mode $\{1\}$.
The restricted Walsh read on one tetrahedral sector therefore has
a natural $3+1$ organisation: three nonconstant spatial characters
and one constant scalar mode.  The same $3+1$ structure appears in
the tetrahedral face decomposition (one uniform-sign face and
three mixed-sign faces~\cite{NelsonA}) and in the Lorentz
signature $(-,+,+,+)$.  Whether these three instances of $3+1$
are independent observations or manifestations of a single
structural fact is left as an open question.

In the companion notes~\cite{NelsonA,NelsonB}, the occupation law
$u(x)=x/(1+x)$ was derived from the sector partition of the
bilinear frame space, and the torsion sector was shown to carry the
regular representation of $S_4$.  The present identification places
both results on the Walsh-Hadamard domain: the sector partition is
a Walsh eigenspace decomposition, and the $S_4$ action is the
permutation group of the Hadamard rows.  The transform that signal
processing uses to analyse Boolean functions is the same transform
that gravity uses to read its frame.


\begin{thebibliography}{99}

\bibitem{Walsh}
J.~L. Walsh,
``A closed set of normal orthogonal functions,''
\emph{Amer.\ J.\ Math.} \textbf{45}, 5--24 (1923).

\bibitem{Sylvester}
J.~J. Sylvester,
``Thoughts on inverse orthogonal matrices\ldots,''
\emph{Phil.\ Mag.} \textbf{34}, 461--475 (1867).

\bibitem{Coxeter}
H.~S.~M. Coxeter,
\emph{Regular Polytopes},
3rd ed., Dover (1973).

\bibitem{Kepler}
J. Kepler,
\emph{Harmonices Mundi}, Linz (1619).
[The stella octangula appears in Book~II.]

\bibitem{MTW}
C.~W. Misner, K.~S. Thorne, and J.~A. Wheeler,
\emph{Gravitation},
W.~H. Freeman (1973).

\bibitem{Hehl1976}
F.~W. Hehl, P.~von der Heyde, G.~D. Kerlick, and J.~M. Nester,
``General relativity with spin and torsion: Foundations and
prospects,''
\emph{Rev.\ Mod.\ Phys.} \textbf{48}, 393--416 (1976).

\bibitem{Hadamard}
J. Hadamard,
``R\'esolution d'une question relative aux d\'eterminants,''
\emph{Bull.\ Sci.\ Math.} \textbf{17}, 240--246 (1893).

\bibitem{NelsonA}
S.~Nelson,
``An algebraic origin for the simple AQUAL interpolation
function,''
submitted (2026).

\bibitem{NelsonB}
S.~Nelson,
``The torsion sector of the tetrad connection as the regular
representation of $S_4$,''
submitted (2026).

\end{thebibliography}

\end{document}
